% OC Core 1.3 -- Modularchitektur und ebeneübergreifende Strukturen

\section{Modularchitektur und ebeneübergreifende Strukturen}
\label{sec:modules_master}

Dieser Abschnitt gibt einen Überblick auf hoher Ebene über die erweiterten strukturellen
Module der Ontologie der Continua. Jedes Modul ist im technischen Atlas ausgearbeitet;
hier fassen wir ihre Rollen und Beziehungen zusammen, sodass die Hauptstruktur eine durchgehende
Brücke zwischen der formalen Doktrin und den Abschlusskapiteln des Core~1.3 bildet.

Die erweiterten Module sind:

\begin{itemize}
\item \textbf{K-levels Master} (Abschnitt~\ref{sec:klevels-master}):
  globaler Überblick über alle Kontinua \(K_0\)–\(K_{12}\), ihre Zustandsräume,
  Achsen, Schwellenwerte und Kernprozesse.

\item \textbf{Einbettungsräume / M-Räume} (Abschnitt~\ref{sec:mspaces-overview}):
  Definition der Meta-Räume \(M_x\), die die Entwicklung der Kontinua \(K_x\)
  einschränken und parametrisieren.

\item \textbf{Cross-level / Cross-K-Strukturen} (Abschnitt~\ref{sec:crossk_master}):
  Operatoren, Abbildungen und Landschaftsstrukturen, die verschiedene Ebenen der
  Hierarchie verbinden und Phasenübergänge zwischen ihnen erzeugen.

\item \textbf{Prozesse} (Abschnitt~\ref{sec:processes-master}):
  vereinheitlichte Beschreibung der Evolutionsoperatoren \(E = (F,G,H,Q,R,S,U)\)
  über Ebenen hinweg, einschließlich Spannungsdynamiken, Schwellen, Flüsse und
  Zyklen.

\item \textbf{Zyklen} (Abschnitt~\ref{sec:cycles-master}):
  strukturelle und dynamische Zyklen \(C_j\) als Träger von Kontinuität,
  Gedächtnis und zeitlicher Organisation.

\item \textbf{Vorhersagen} (Abschnitt~\ref{sec:predictions-master}):
  ebene-spezifische und ebeneübergreifende Vorhersagen, empirische Signaturen und
  qualitative Regime, die durch die Formalisierung impliziert werden.

\item \textbf{Experimente} (Abschnitt~\ref{sec:experiments-master}):
  konzeptionelle und konkrete experimentelle Designs zur Erprobung von Kontinua,
  Übergängen und Schwellenlandschaften.

\item \textbf{Falsifizierbarkeitsstruktur} (Abschnitt~\ref{sec:falsifiability-master}):
  explizite Bedingungen, unter denen die Ontologie der Continua widerlegt,
  eingegrenzt oder zu einer Neuordnung ihres Kerns gezwungen werden kann.

\item \textbf{Jets / lokale Evolutionsstruktur} (Abschnitt~\ref{sec:jets-master}):
  lokale Entwicklungen des Evolutionsoperators und Kurzzeitdynamiken rund um
  gegebene Konfigurationen von \(K_x\) und \(M_x\).
\end{itemize}

Gemeinsam verfeinern diese Module die Kernformalismen über Ebenen und Anwendungen
hinweg, wobei sie zugleich eine einzige kohärente Hierarchie von Kontinua und
Metaräumen erhalten.
